\section{Partition-Function Structure}
\label{sec:partition}

\subsection{The Cumulant-Generating Identity}

The canonical form \cref{eq:canonical} yields an immediate log-identity:
\begin{equation}
  \log\eta(\tau,\rho)
  = -\log\!\bigl(1 - \rho\,\sech\tau\bigr).
  \label{eq:logeta_identity}
\end{equation}
Setting $u = \rho\,\sech\tau$ and $\theta = \log u$ for $u \in (0,1)$, the right-hand side becomes
\[
  -\log(1 - e^{\theta}) = \sum_{n=1}^{\infty}\frac{e^{n\theta}}{n},
\]
which is the log-partition function (cumulant generating function at zero argument) of the \emph{logarithmic series distribution}. A more elementary identification comes from the geometric series:
\[
  \eta = \frac{1}{1 - u} = \sum_{n=0}^{\infty}u^{n},\qquad u \in (0,1),
\]
recognising $u^n = e^{n\theta}$ as the natural-parameter exponential-family weighting of a latent count $n \in \{0,1,2,\dotsc\}$.

\begin{theorem}[Geometric Exponential Family]
\label{thm:exp_family}
For $\rho\,\sech\tau \in (0,1)$, define the natural parameter $\theta = \log(\rho\,\sech\tau) < 0$ and the log-partition function $A(\theta) = \log\eta$.  Then
\begin{equation}
  p(n;\theta) = (1-e^{\theta})\,e^{n\theta},\qquad n \in \{0,1,2,\dotsc\},
  \label{eq:geo_family}
\end{equation}
defines a one-parameter exponential family with sufficient statistic $N$, base measure $h(n) = 1$, and $A(\theta) = -\log(1-e^{\theta}) = \log\eta$.  The PEF $\eta(\tau,\rho)$ is the partition function of this family, evaluated at the effective parameter $\theta = \log(\rho\,\sech\tau)$ determined by the pair $(\tau,\rho)$.
\end{theorem}

\begin{proof}
Normalisation: $\sum_{n=0}^{\infty}(1-e^{\theta})e^{n\theta} = (1-e^{\theta})/(1-e^{\theta}) = 1$ for $e^{\theta} \in (0,1)$.  The cumulant generating function is $\log\mathbb{E}[e^{tN}] = -\log(1-e^{t+\theta}) + \log(1-e^{\theta})$, so $A(\theta) = -\log(1-e^{\theta}) = \log\eta$ is the log-partition function evaluated at $t=0$.
\end{proof}

\subsection{First and Second Cumulants}
\label{sec:cumulants}

The standard exponential-family relations $\mathbb{E}[N] = A'(\theta)$ and $\Var(N) = A''(\theta)$ give:

\begin{proposition}[Cumulant Identities]
\label{prop:cumulants}
\[
  \mathbb{E}[N] = A'(\theta) = \frac{e^{\theta}}{1-e^{\theta}} = \eta - 1,
  \qquad
  \Var(N) = A''(\theta) = \frac{e^{\theta}}{(1-e^{\theta})^{2}} = \eta(\eta-1).
\]
\end{proposition}

\begin{proof}
$A'(\theta) = e^{\theta}/(1-e^{\theta})$.  Since $e^{\theta} = u = \rho\,\sech\tau$ and $1-u = 1/\eta$, we get $A'(\theta) = u\eta = \eta - 1/\eta \cdot \eta = \eta - 1$.  For the second derivative, $A''(\theta) = e^{\theta}/(1-e^{\theta})^{2} = u\eta^{2} = (\eta-1)\eta$.
\end{proof}

The identities admit a direct empirical reading.  The quantity $\eta - 1$ is the \emph{signal}: for $\eta > 1$ it is positive, measuring the fractional gain from pairing; for $\eta < 1$ it is negative, measuring the fractional loss.  The quantity $\eta(\eta-1)$ is the \emph{curvature of the log-partition function}, which determines how rapidly $A$ changes with the natural parameter: it is the Fisher information of the latent family, and it sets the scale of fluctuations of $\hat{\eta}$ around its expectation.  Both quantities are properties of $\eta$ alone, valid for any $(\kappa,\rho)$ such that $\rho\,\sech\tau \in (0,1)$, with no additional distributional assumptions beyond the PEF definition.

\subsection{Regime Change at \texorpdfstring{$\rho = 0$}{rho = 0}}
\label{sec:regime}

The geometric-family identification requires $\rho\,\sech\tau \in (0,1)$, i.e.\ $\rho > 0$.  For $\rho < 0$, the series $\sum_{n=0}^{\infty}(\rho\,\sech\tau)^{n}$ becomes sign-alternating: even-$n$ terms contribute positively, odd-$n$ terms negatively.  The closed form $\eta = 1/(1-\rho\,\sech\tau)$ remains well-defined and finite (since $|\rho\,\sech\tau| < 1$ in the physical domain), but the partition-function reading breaks down because $\theta = \log(\rho\,\sech\tau)$ is undefined for $\rho < 0$.

There is a symmetry that partially restores the picture: write $\eta(\tau,\rho)$ for $\rho < 0$ as
\[
  \eta(\tau,\rho)
  = \frac{\cosh\tau}{\cosh\tau + |\rho|}
  = 1 - \frac{|\rho|}{\cosh\tau + |\rho|} < 1.
\]
Setting $\tilde{u} = |\rho|\,\sech\tau \in (0,1)$, we have $\eta = 1 - \tilde{u}/(1+\tilde{u}) = 1/(1+\tilde{u})$, which is distinct from the $\rho > 0$ form $1/(1-u)$.  The two branches meet at $\eta = 1$ when $\rho = 0$.  The partition function structure is therefore intrinsically chiral: it requires a sign choice at $\rho = 0$ before the latent-geometric reading applies.

This is the mathematical origin of the regime change that the empirical paper terms the `efficiency--power tension' in competitive settings.  The prediction is concrete: statistics and machine-learning diagnostics that are smooth functions of $\eta$ will show a detectable slope discontinuity at $\rho = 0$, reflecting the change in the curvature $A''(\theta)$ across the boundary.  \Cref{sec:val_regime} tests this empirically.

\subsection{Distribution-Free Scope}
\label{sec:distrib_free}

The Gaussian mutual-information link of the empirical companion \parencite{brownPEFempirical} requires bivariate normality.  The cumulant identities of \cref{sec:cumulants} do not: they follow from the algebraic form of $\eta$ alone.  Specifically, \cref{prop:cumulants} uses no properties of the original data-generating distribution beyond the definition of $\kappa$ and $\rho$ as (population) variance ratio and correlation.  The result is:
\begin{itemize}
  \item the `signal' quantity $\eta-1$ and the `information' quantity $\eta(\eta-1)$ are \emph{distribution-free} properties of the efficiency landscape;
  \item the Gaussian mutual information $I(X;Y) = -\tfrac{1}{2}\log(1-\rho^{2})$ is a $(\kappa,\rho)$-specific instance arising only when $\kappa = 1$ (Fisher's case) under joint normality;
  \item the partition-function form $\log\eta = -\log(1-\rho\,\sech\tau)$ is the natural generalisation that replaces the Gaussian special case with a formula that is valid across all variance ratios and correlations without distributional assumptions.
\end{itemize}

\subsection{The Latent Count}
\label{sec:latent_count}

The random variable $N$ in \cref{thm:exp_family} has no direct interpretation as a counting variable in the original $(X_{\mathrm{A}},X_{\mathrm{B}})$ problem.  It is a \emph{latent} sufficient statistic in the sense of exponential-family theory \parencite{brown1986,coverThomas2006}: a formal object whose existence is implied by the algebraic form of $A(\theta)$, not by any physical counting process.  Analogous latent constructions appear throughout statistics; the Poisson--Binomial relationship and the Pólya-urn representation of the Dirichlet--Multinomial are comparable examples.

The interpretation we adopt is functional rather than generative: $\mathbb{E}[N] = \eta - 1$ quantifies deviation from the pairing-null ($\eta = 1$, $N \equiv 0$), and $\Var(N) = \eta(\eta-1)$ quantifies the sensitivity of the log-partition function to changes in the effective parameter $\theta = \log(\rho\,\sech\tau)$.  These are empirically accessible through the estimator $\hat{\eta}$, and their validity is independent of the data-generating mechanism.
